\section{Related Work}\label{sec:related}

\noindent\textbf{Local Planners and Topological Consistency.} \texttt{DWA}~\citep{fox1997dwa} selects a velocity command from a dynamic window at every cycle; each cycle's selection is made anew, with no memory of previous decisions, so nothing in the scheme itself discourages the chosen side from flipping between cycles (a structural observation of ours, not a claim made by \citet{fox1997dwa}). \texttt{TEB}~\citep{rosmann2017teb} maintains and optimizes several topological alternatives in parallel; its ROS implementation adds a selection-cost hysteresis that discourages switching between alternatives, but this is an implementation-level heuristic --- the formulation itself does not define commitment across cycles. \texttt{MPPI}~\citep{williams2017mppi} is a sampling-based MPC formulation; it does not treat the discrete structure of social passing decisions as first class. H-signature-based enumeration of homotopy classes~\citep{bhattacharya2012topological} rigorously represents the topological identity of a path but focuses mainly on static-obstacle topology; the branch that extends this topological view to moving-pedestrian environments is most relevant to this paper. Topological braids model the joint avoidance strategy of several agents, and a planner can then comply with the avoidance intentions inferred from their past behavior~\citep{mavrogiannis2019braids}. Dynamic Channel~\citep{cao2019dynamicchannel} is a real-time crowd-navigation framework that searches a triangulated spatial graph, combining pedestrian dynamics with topological-corridor selection; it picks one topology per planning cycle but does not formalize the temporal consistency of that choice via an explicit commitment rule or switching-rate guarantee. Winding Through~\citep{mavrogiannis2023winding} tracks the homotopy class the robot's path holds relative to each pedestrian and penalizes changes to that class over time as a cost --- being motivated by suppressing passing-side switches, it is among the closest prior work, but it operates as a soft topological-invariance cost internal to the planner and does not define a track-ID-indexed \Class\ lifecycle or a provable switching-rate bound. Topology-driven parallel trajectory optimization~\citep{degroot2024topology} runs local MPC optimizers seeded in different homotopy classes in parallel for dynamic obstacles and commits to the best resulting trajectory each cycle, maintaining consistency of the chosen class --- this treats homotopy consistency via parallel search at the trajectory-optimization layer, and does not provide an explicit accumulator-based commitment rule at the discrete decision layer nor a formal switching-rate guarantee for it. SHINE~\citep{martinezbaselga2026shine} selects among topologically distinct guidance trajectories with a network trained on human motion data and reduces the selection cost of the homology class followed in the previous control period by a consistency weight, stating explicitly that fast switching between topology classes can lead to indecisive and ultimately non-social navigation. SHINE therefore treats decision oscillation in social navigation as a design problem; its remedy is a one-step, previous-class cost bias, and it does not state a bound on switching frequency. \citet{zhou2023minimally} keep the previously selected trajectory as a candidate with a time-decaying weight, motivated by the observation that a high decision frequency, together with perception noise, can keep pedestrians from discerning a stable intention and cause oscillation. In multi-agent navigation, a learned planner can choose target winding numbers that a model-predictive controller then realizes~\citep{nakao2025wnummpc}, a decision--control split aimed at breaking symmetry deadlocks rather than at bounding switching. \sname\ differs from these mainly in making the cross-cycle commitment rule itself, rather than a per-cycle cost bias, the object of analysis, and in specifying \Class\ correspondence over time (lifecycle) against track identity.

\noindent\textbf{Explicit Passing-Side Decisions and Opinion Dynamics.} Opinion-dynamics navigation makes the passing side an explicit decision state. \citet{cathcart2023opinion} drive human--robot corridor passing with nonlinear opinion dynamics in which the robot forms an opinion about which side to pass an approaching person; they show analytically that deadlock in the decision is broken when the robot's attention exceeds a critical value, even when the robot has no bias or evidence for either side, and that opinions form rapidly; they validate these predictions in human--robot experiments. Opinion dynamics has also been combined with vortex fields so that robot and humans reach a consensus on the direction of motion~\citep{daddato2024oval}, and with dynamical-system modulation so that the robot chooses a passing side from onboard sensing, with simulation and human-participant experiments~\citep{chen2025socialdsm}. These methods are the closest prior work on a stateful, analyzed decision layer for the passing side. The formal target of \citet{cathcart2023opinion} is the formation of a decision (deadlock breaking); that of \sname\ is a bound on how often a formed decision is revised, over multi-person track-indexed labels. \sname\ has no comparable human or closed-loop evidence. Concurrently with this work, BIG-CBF~\citep{li2026bigcbf} separates maneuver selection from safety enforcement: it selects among six closed-loop behaviors, including passing left and right, with a switching cost relative to the retained mode, a minimum dwell time and a relative switching-improvement threshold, gives a hard control-barrier-function filter final safety authority, and reports closed-loop simulation and physical-robot experiments. BIG-CBF thus covers mode retention, dwell and improvement gating, and safety priority, with much broader empirical evidence than this paper; it does not use an evidence accumulator, a progress-dependent threshold or track-indexed topological classes, and we did not find a stated switching-frequency bound in it.

\noindent\textbf{Change Detection, Switching Logic, and Arbitration.} The mathematics used in this paper is standard. Page's cumulative-sum (CUSUM) inspection scheme~\citep{page1954cusum} accumulates the deviations of observations from a reference value and signals when the cumulative sum rises by a decision interval above its running minimum; in recursive form this is a sum reflected at zero. With the leak removed and a constant threshold, the accumulator of \cref{subsec:switch} is this statistic, written as Lindley's queueing recursion~\citep{lindley1952queue}, and the exponential bound of P1 is the classical tail bound for that recursion~\citep{kingman1970queues}. Leaky accumulation of evidence toward a threshold is a standard model of perceptual choice~\citep{usher2001lca}. In supervisory control, hysteresis switching logic changes the active candidate only when its monitoring signal reaches $(1+h)$ times the smallest one, for a hysteresis constant $h>0$, and yields a bound on the number of switches on an arbitrary time interval~\citep{hespanha2003hysteresis}; dwell-time switching separates consecutive switches by at least a prescribed interval, and an average dwell-time bounds the number of switches on any interval by a chatter bound plus a term linear in the interval length~\citep{hespanha1999dwell}. P2 is a dwell-time guarantee of this kind for the discretionary branch, obtained by an elementary bounded-increment argument. Finally, placing a safety behavior strictly above nominal behaviors is standard priority-based behavior arbitration~\citep{orzechowski2020arbitration}, which has been extended with verification and fallback layers~\citep{spieker2024arbitration}. We use these results as tools and claim neither the accumulator nor its bounds as new mathematics.

\noindent\textbf{Human Models, Prediction, and Learned Social Costs.} The Social Force Model~\citep{helbing1995social} describes pedestrian motion as interaction forces, and \texttt{ORCA}/\texttt{RVO}~\citep{vandenberg2011rvo} formalizes reciprocal multi-agent collision avoidance; social-navigation work that addresses freezing models such human reactivity and cooperation explicitly~\citep{trautman2010,trautman2015dense}. Proxemics~\citep{hall1966hidden} contributes qualitative, culture-dependent interpersonal distance zones, which robotics work has operationalized as cost fields --- notably the asymmetric-Gaussian formulation of \citet{kirby2010thesis}. Learning-based probabilistic predictors --- Social-LSTM~\citep{alahi2016sociallstm}, Social-GAN~\citep{gupta2018socialgan}, Trajectron++~\citep{salzmann2020trajectron} --- predict a person's future trajectory as a distribution, and the line of work that learns socially compliant costs and prediction models from pedestrian trajectory features via maximum-entropy inverse reinforcement learning~\citep{kuderer2012feature,kretzschmar2016irl} models human cooperative avoidance behavior as a trajectory distribution (including a mixture distribution over joint path choices), generating robot trajectories that conform to observed walking conventions. While these provide rich human models and learned cost/prediction models for trajectory generation, the works cited in this paragraph do not treat the oscillation of the robot's decision itself as a measurement target or design goal, nor state switching-rate bounds for a discrete passing decision. Learned topology selection differs in this respect: SHINE~\citep{martinezbaselga2026shine} explicitly treats fast class switching as a design problem (first paragraph of this section). \sname\ adopts such cost fields as input and uses a Kalman constant-velocity prediction over a short horizon (\cref{subsec:formulation,sec:disc}); the learned predictors can be incorporated as a replaceable input module, since this contribution is orthogonal to the choice of predictor --- only the mean-trajectory and covariance input need be replaced.

\noindent\textbf{Social Momentum, Legibility, and Group Awareness.} Cooperation-aware navigation that jointly models robot and human trajectories in planning~\citep{trautman2015dense} addresses cooperative navigation amid dense crowds. In particular, social momentum~\citep{mavrogiannis2022momentum} scores candidate actions by the angular momentum of the joint robot--human configuration to reconcile human--robot passing-side intentions and reinforce implicit coordination; it is a per-cycle objective rather than a persistent accumulator, but it articulates the intuition --- committing to a passing side aids coordination --- that \sname\ formalizes as an explicit scalar evidence accumulator. Legible/predictable motion~\citep{dragan2013legibility} presents the view that robot motion should reveal intent to an observer; the commitment mechanism of \sname\ suppresses oscillation of the discrete decision label, which we hypothesize can support legible motion; this paper does not show it, because the offline observer scores decision-label streams only (\cref{subsec:setup}) and motion legibility is not measured. The methodology for human validation of legibility is elaborated via the observer model of \cref{subsec:observer} together with the user-study protocol of \cref{app:planned}. Metrics for human comfort and social appropriateness are undergoing standardization in social-navigation surveys~\citep{mavrogiannis2023survey,gao2022evaluation,francis2023principles}, and \sname\ takes the axes of decision oscillation and legibility as its quantitative target among these. The spatial arrangement people form during conversation/interaction (F-formation) was formalized by \citet{kendon1990conducting} via the concepts of o-space and p-space; treating such formations as ``groups that must not be split'' is a norm adopted in robot navigation, not part of Kendon's own definition, and detecting and respecting F-formations has been treated as a core element of group-aware navigation~\citep{riosmartinez2015survey}. Automatic F-formation detection itself is outside the scope of this paper and is assumed to be provided by an external perception module; \sname\ extends the merge rule of \cref{subsec:enum} to interaction-based groups so as to avoid splitting F-formations (\cref{subsec:groups}).

\noindent\textbf{Positioning.} The problem of temporal consistency on the passing side has been recognized in prior work, with mechanisms of varying explicitness --- a soft topological-invariance penalty~\citep{mavrogiannis2023winding}, the selection-cost hysteresis in the ROS implementation of \texttt{TEB}~\citep{rosmann2017teb}, per-cycle topological-corridor selection~\citep{cao2019dynamicchannel} (one corridor is chosen each cycle, but no cross-cycle commitment is defined), the social-momentum objective~\citep{mavrogiannis2022momentum}, and previous-trajectory retention~\citep{zhou2023minimally}. \citet{degroot2024topology} \emph{do} formalize inter-cycle topological consistency: a consistency parameter trades off following the previously selected homotopy class against switching to a better one, and at one extreme enforces the previous class as long as it remains valid; SHINE applies a previous-class consistency weight in social navigation for the stated purpose of avoiding indecisive navigation~\citep{martinezbaselga2026shine}. Opinion-dynamics methods make the passing side an explicit, stateful decision with analytical deadlock-breaking guarantees~\citep{cathcart2023opinion}, and concurrent work combines passing-mode retention, dwell and relative-improvement gating, and a hard safety filter in closed loop~\citep{li2026bigcbf}. The mechanisms and mathematics we use are also standard (see Change Detection, Switching Logic, and Arbitration above). The delta of this paper is therefore narrow, and we state it precisely: (i) a leaky, CUSUM-type evidence accumulator with a progress-dependent threshold, used as the commitment rule for a social passing decision; (ii) an explicit switching-frequency contract for that decision --- the deterministic, pathwise inter-switch bound P2, its conditional exponential refinement P1, and the conditional response-time bound --- stated together with the assumptions each needs; (iii) a track-ID-indexed \Class\ representation with lifecycle rules (spawn, removal, merge, split, TTL), which defines consistency against person identity rather than against a geometric homotopy class recomputed from the current scene, specified as a design contract that the current implementation only partly realizes; and (iv) an open ROS~2 Nav2 controller plugin with an offline harness whose ablations indicate that, in the tested regime, the accumulator rather than the margin or progress hardening carries the oscillation suppression. To the author's knowledge, no prior social-navigation work states a switching-frequency bound for the discrete passing decision. We claim neither the individual mechanisms nor the mathematics as new, and we have not compared against consistency-weight, opinion-dynamics or dwell-plus-hysteresis baselines (\cref{sec:disc}).
